% Code Listings
% Demonstrates: listings package for syntax-highlighted code in multiple languages.
% Compile with: pdflatex main.tex
% Note: For minted (Pygments-based), compile with: pdflatex -shell-escape main.tex

\documentclass[12pt]{article}
\usepackage[utf8]{inputenc}
\usepackage[T1]{fontenc}
\usepackage{xcolor}
\usepackage{listings}
\usepackage{hyperref}

% =============================================================================
% Listings configuration
% =============================================================================

% --- Define a color scheme ---
\definecolor{codebg}{RGB}{245,245,245}
\definecolor{codegreen}{RGB}{0,128,0}
\definecolor{codegray}{RGB}{128,128,128}
\definecolor{codepurple}{RGB}{128,0,128}
\definecolor{codeblue}{RGB}{0,0,180}

% --- Default style applied to all listings ---
\lstset{
    backgroundcolor=\color{codebg},
    basicstyle=\ttfamily\small,
    breaklines=true,
    captionpos=b,
    commentstyle=\color{codegreen}\itshape,
    frame=single,
    keywordstyle=\color{codeblue}\bfseries,
    numbers=left,
    numberstyle=\tiny\color{codegray},
    numbersep=8pt,
    showstringspaces=false,
    stringstyle=\color{codepurple},
    tabsize=4,
    xleftmargin=1.5em,
    framexleftmargin=1em,
}

\title{Code Listings in \LaTeX{}}
\author{LaTeX Tutorial}
\date{\today}

\begin{document}
\maketitle

% =============================================================================
\section{Python}
% =============================================================================

\begin{lstlisting}[language=Python, caption={Fibonacci with memoization}]
def fibonacci(n, memo={}):
    """Return the n-th Fibonacci number."""
    if n in memo:
        return memo[n]
    if n <= 1:
        return n
    memo[n] = fibonacci(n - 1) + fibonacci(n - 2)
    return memo[n]

for i in range(10):
    print(f"F({i}) = {fibonacci(i)}")
\end{lstlisting}

% =============================================================================
\section{C}
% =============================================================================

\begin{lstlisting}[language=C, caption={Hello World in C}]
#include <stdio.h>

int main(void) {
    printf("Hello, World!\n");
    return 0;
}
\end{lstlisting}

% =============================================================================
\section{Java}
% =============================================================================

\begin{lstlisting}[language=Java, caption={A simple generic stack}]
import java.util.ArrayList;

public class Stack<T> {
    private final ArrayList<T> items = new ArrayList<>();

    public void push(T item) {
        items.add(item);
    }

    public T pop() {
        if (items.isEmpty()) {
            throw new RuntimeException("Stack is empty");
        }
        return items.remove(items.size() - 1);
    }

    public boolean isEmpty() {
        return items.isEmpty();
    }
}
\end{lstlisting}

% =============================================================================
\section{Inline Code}
% =============================================================================

Use \lstinline|int x = 42;| for inline code, or
\lstinline[language=Python]|print("hello")| with a language override.

% =============================================================================
\section{Custom Language Definition}
% =============================================================================

You can define your own language for DSLs or config files:

\lstdefinelanguage{TOML}{
    comment=[l]{\#},
    morestring=[b]",
    morekeywords={true, false},
}

\begin{lstlisting}[language=TOML, caption={Example TOML configuration}]
# Application settings
[server]
host = "0.0.0.0"
port = 8080
debug = false

[database]
url = "postgres://localhost/mydb"
pool_size = 5
\end{lstlisting}

% =============================================================================
\section{Loading from External File}
% =============================================================================

To include code from an external file, use:
\begin{verbatim}
\lstinputlisting[language=Python, caption={My script}]{script.py}
\end{verbatim}

This keeps your \LaTeX{} source clean and your code in separate files.

\end{document}
